\documentclass[11pt,a4paper]{article}
\usepackage[T1]{fontenc}
\usepackage{lmodern}
\usepackage[top=0.45in, bottom=0.45in, left=0.68in, right=0.68in, headheight=0pt, headsep=0pt]{geometry}
\usepackage{amsmath,amssymb}
\usepackage{xcolor}
\usepackage{tikz}
\usetikzlibrary{arrows.meta,positioning,calc,decorations.pathreplacing}
\usepackage{tcolorbox}
\usepackage{hyperref}
\usepackage{booktabs}
\usepackage{caption}

\hypersetup{
    colorlinks=true,
    linkcolor=blue!70!black,
    citecolor=blue!70!black,
    urlcolor=blue!70!black
}

% Custom Callout Boxes matching math_tools_en / house style
\newtcolorbox{learningobj}{
    colback=blue!3!white,
    colframe=blue!60!black,
    fonttitle=\bfseries\sffamily,
    fontupper=\small,
    title={$\blacktriangleright$ LEARNING OBJECTIVE},
    arc=1.5mm,
    boxrule=0.8pt,
    left=3mm, right=3mm, top=1.2mm, bottom=1.2mm,
    before skip=1.0mm, after skip=2.0mm
}

\newtcolorbox{groundbreaking}{
    colback=teal!4!white,
    colframe=teal!60!black,
    fonttitle=\bfseries\sffamily,
    fontupper=\small,
    title={$\bigstar$ GROUNDBREAKING FINDING},
    arc=1.5mm,
    boxrule=0.8pt,
    left=3mm, right=3mm, top=1.2mm, bottom=1.2mm,
    before skip=1.5mm, after skip=2.0mm
}

\newtcolorbox{physicalinsight}{
    colback=yellow!6!white,
    colframe=orange!75!black,
    fonttitle=\bfseries\sffamily,
    fontupper=\small,
    title={$\blacktriangleright$ PHYSICAL INSIGHT},
    arc=1.5mm,
    boxrule=0.8pt,
    left=3mm, right=3mm, top=1.2mm, bottom=1.2mm,
    before skip=1.5mm, after skip=2.0mm
}

\renewcommand{\thefigure}{1.4\alph{figure}}

\begin{document}

% ============================================================================
% PAGE 1: Learning Objective, Bridge, State 1, Table 1.4, Figure 1.4a
% ============================================================================
\vspace*{-4.5mm}
\section*{1.4 Hardware Consequences and Processing Latency Limits}
\vspace{-2.5mm}

\begin{learningobj}
Rather than treating hardware latency as an arbitrary processing delay, this section grounds timing in physical acoustic accumulation, contrasts on-chip Block RAM allocations against external memory bandwidth, and exposes why the migration of audio front-ends to FPGA logic is driven by deterministic execution and zero-jitter hardware isolation rather than raw arithmetic throughput.
\end{learningobj}

\vspace{0.5mm}
\noindent\textbf{Organic Bridge from Section 1.3: Mapping Equations onto Physical Transistors.} In Section 1.3, we completed the mathematical microscope of the front-end DSP pipeline, transforming 400 streaming pressure samples into 80 Log-Mel spectral values. But an algorithm on paper assumes infinite memory and zero execution time. In this section, we force these abstract equations to collide with physical silicon: evaluating how the mathematical parameters translate into memory footprint and processing deadlines across an edge GPU and a spatial FPGA fabric (AMD Xilinx Kria KV260).

\vspace{-1.5mm}
\subsection*{State 1: The Physical Latency Floor}
\vspace{-2mm}

Latency in a streaming audio system has an absolute physical lower bound dictated by acoustics before any semiconductor chip is ever powered on. Three temporal quantities are permanently locked by the front-end acoustic specification:

\vspace{-2mm}
\begin{center}
\small
\begin{tabular}{l c p{10.2cm}}
\toprule
\textbf{Parameter} & \textbf{Value} & \textbf{Mathematical Origin \& Physical Definition} \\
\midrule
First-frame accumulation floor & $25\text{ ms}$ & $L / f_s = 400 / 16{,}000$ (Microphone diaphragm collection time) \\
Hop cadence between frames & $10\text{ ms}$ & $H / f_s = 160 / 16{,}000$ (Periodic emission interval to downstream model) \\
Overlap ratio across frames & $60\%$ & $(L - H) / L = 240 / 400$ (Shared historical context preserving energy) \\
\bottomrule
\end{tabular}
\end{center}
\vspace{-3.5mm}

This acoustic reality establishes a strict latency floor that no processor can circumvent: \textbf{Frame 0 physically cannot exist} before $25\text{ ms}$ of continuous air pressure vibrations have physically struck the microphone diaphragm. No processor, regardless of clock frequency, can circumvent this acoustic floor: Frame 0 cannot exist before $25\text{ ms}$ because sample 400 has not arrived.

\vspace{-2.5mm}
\begin{center}
\begin{minipage}{\textwidth}
\centering
\resizebox{0.94\textwidth}{!}{\input{chapter01/fig_1_4a.tex}}\\[0.8mm]
\begin{minipage}{0.96\textwidth}
\small\textbf{Figure 1.4a}: \textbf{The acoustic latency floor versus processor compute speed.} (1)~\textbf{What it is:} Streaming audio timeline showing physical frame accumulation and periodic hop emission. (2)~\textbf{What to look at:} Frame 0 requires a full $25\text{ ms}$ ($400$ samples) of acoustic pressure accumulation before processing can begin; Frame 0 cannot exist before $25\text{ ms}$ because sample 400 has not arrived. Subsequent frames advance every $10\text{ ms}$ ($160$ samples) with a $60\%$ ($240$-sample) overlap. (3)~\textbf{What it proves:} Real-time latency in voice edge systems is strictly governed by physical wave propagation and windowing acoustics, not processor clock frequency.
\end{minipage}
\end{minipage}
\end{center}
\vspace{-3.5mm}

Once the initial $25\text{ ms}$ accumulation floor is satisfied, the system enters steady-state operation. From this point forward, the microphone emits a new block of $H = 160$ samples every $10\text{ ms}$. This $10\text{ ms}$ interval is the \textbf{hard real-time deadline}: all downstream computations---windowing, 512-point FFT, power spectrum, 80-channel Mel filtering, logarithmic compression, and neural network inference---must complete within $10\text{ ms}$. A faster processor does not shrink the $25\text{ ms}$ physical floor; it merely prevents additional computational delay from accumulating on top of it.

\clearpage

% ============================================================================
% PAGE 2: State 2, Physical Insight, State 3, Throughput vs Determinism
% ============================================================================
\vspace*{-4.5mm}
\subsection*{State 2: Memory Footprint on Physical Silicon}
\vspace{-2mm}

To evaluate physical feasibility, we ground our memory requirements against official hardware specifications from the AMD DS890 datasheet for the Zynq UltraScale+ \texttt{XCK26} MPSoC deployed on the Kria KV260:
\textbf{117,120} logic cell Lookup Tables (LUTs); \textbf{144} Block RAM blocks of $36\text{ kbit}$ each ($5.1\text{ Mb} \approx 648\text{ KiB}$ total on-chip SRAM); \textbf{64} UltraRAM blocks ($18.0\text{ Mb} \approx 2.25\text{ MiB}$ total bulk SRAM); and \textbf{1,248} DSP48E2 compute slices optimized for single-cycle multiply-accumulates.

Now contrast these available resources against the two core state structures that the streaming DSP front-end must maintain in 32-bit single-precision floating point (FP32):
\begin{itemize}\setlength{\itemsep}{0.5mm}\setlength{\parskip}{0pt}
  \item \textbf{Audio Ring Buffer ($400$ samples)}: At $4\text{ bytes}$ per sample, storing running acoustic history requires $400 \times 4\text{ B} = 1{,}600\text{ bytes}$. A single $36\text{-kbit}$ Block RAM block provides $4{,}608\text{ bytes}$ of dual-port storage. The entire time-domain state occupies only \textbf{$34.7\%$ of a single BRAM block} ($0.2\%$ of total on-chip BRAM).
  \item \textbf{Mel Filterbank Weight Matrix ($80 \times 257$)}: Dense storage of triangular filter weights requires $80 \times 257 \times 4\text{ B} = 82{,}240\text{ bytes} = 80.3\text{ KiB}$. This entire coefficient matrix requires 18 Block RAM blocks, consuming exactly \textbf{$12.4\%$ of total on-chip BRAM}.
\end{itemize}

\begin{physicalinsight}
Because the combined memory footprint of the front-end state and filterbank matrix ($81.9\text{ KiB}$) consumes less than $13\%$ of on-chip Block RAM, the entire DSP pipeline fits within local FPGA SRAM. The hardware requires zero round-trips to external DDR DRAM, eliminating dynamic memory arbitration, DRAM refresh stalls, and memory bus energy consumption.
\end{physicalinsight}

\vspace{-2.5mm}
\subsection*{State 3: The Arithmetic Fallacy --- Throughput vs.\ Deterministic Isolation}
\vspace{-2mm}

Evaluating the 80-channel Mel filterbank requires multiplying 257 power spectral bins by 80 triangular weight vectors:
$$\text{Workload per frame} = 80 \times 257 = 20{,}560\ \text{MAC operations}$$
At a hop cadence of $10\text{ ms}$, the front-end produces $100\text{ frames per second}$, resulting in an aggregate computational throughput of:
$$\text{Throughput} = 20{,}560\ \text{MAC/frame} \times 100\ \text{frames/s} = 2.056 \times 10^6\ \text{MAC/s} \approx 2.06\ \text{MMAC/s}$$

Spread evenly across the $1{,}248$ DSP48E2 slices of the Kria KV260, each individual DSP slice would need to execute only \textbf{$1{,}647\text{ operations per second}$}. When clocked at $200\text{ MHz}$, a single DSP48E2 slice provides $200 \times 10^6\text{ MAC/s}$; thus, the entire front-end arithmetic workload represents less than $0.001\%$ of the chip's computational capacity.

Now compare this to the competing edge accelerator: a commercial edge GPU accelerator. A modern edge GPU delivers up to tens of sparse INT8 TOPS at a $25\text{ W}$ power ceiling. Against this massive computing budget, the $2.06\text{ MMAC/s}$ required by the voice front-end is literally \textbf{one millionth} of the GPU's available processing bandwidth.

This disparity exposes a central architectural truth: \textbf{migrating an audio front-end to FPGA logic is never about raw arithmetic throughput}. An edge GPU possesses orders of magnitude more raw floating-point performance. The migration is driven entirely by \textbf{temporal determinism and hardware isolation}:
\begin{itemize}\setlength{\itemsep}{0.5mm}\setlength{\parskip}{0pt}
  \item \textbf{The Shared Software Trap (Edge GPU)}: On a GPU/CPU system, the audio front-end shares the $10\text{ ms}$ processing window with the Linux operating system kernel, microphone DMA driver interrupts, background network stacks, and concurrent neural network inference. Each CUDA kernel invocation incurs $5\text{--}15\ \mu\text{s}$ of launch overhead, and operating system thread scheduling introduces random preemption. A deadline shared among multiple asynchronous software layers is inherently subject to latency jitter ($\Delta t_{\text{jitter}}$).
  \item \textbf{The Dedicated Spatial Fabric (FPGA)}: On an FPGA, the DSP pipeline is constructed from dedicated, physically isolated silicon wires and registers. It does not share a clock, an execution pipeline, or a memory bus with an operating system. The front-end completes each frame in a mathematically fixed number of clock cycles ($N_{\text{clk}}$), guaranteeing an execution jitter of identically zero nanoseconds ($\Delta t = 0\text{ ns}$).
\end{itemize}

\clearpage

% ============================================================================
% PAGE 3: Figure 1.4b, Groundbreaking Finding, Organic Bridge to Section 1.5
% ============================================================================
\vspace*{-4.5mm}
\begin{center}
\begin{minipage}{\textwidth}
\centering
\resizebox{0.94\textwidth}{!}{\input{chapter01/fig_1_4b.tex}}\\[1mm]
\begin{minipage}{0.96\textwidth}
\small\textbf{Figure 1.4b}: \textbf{KV260 BRAM memory breakdown and datapath determinism.} (1)~\textbf{What it is:} Block RAM resource distribution on the AMD XCK26 MPSoC alongside architectural pipeline comparison. (2)~\textbf{What to look at:} In (a), front-end state ($1{,}600\text{ B}$) and Mel weights ($80.3\text{ KiB}$) occupy only $12.6\%$ of on-chip SRAM, leaving $87.4\%$ for acoustic neural networks. In (b), edge GPU shares the $10\text{ ms}$ hop window with OS kernel interrupts and memory contention (jitter), whereas the FPGA datapath executes with dedicated silicon wires and zero jitter. (3)~\textbf{What it proves:} The migration imperative from GPU to FPGA is governed by deterministic latency isolation, not raw FLOPS.
\end{minipage}
\end{minipage}
\end{center}
\vspace{-2.5mm}

\begin{groundbreaking}
The FPGA does not outperform the edge GPU on raw arithmetic throughput; it outperforms the GPU on deterministic isolation. Dedicated silicon wires guarantee zero-jitter execution within the non-negotiable 10 ms hop deadline, completely decoupled from operating system preemption and memory contention.
\end{groundbreaking}

\vspace{2mm}
\noindent\textbf{Organic Bridge to Section 1.5: Stress-Testing Hardware Boundaries.} We have established that the streaming front-end pipeline maps with exceptional efficiency onto physical FPGA resources---consuming minimal Block RAM and a negligible fraction of DSP slices while securing absolute latency determinism. 

However, mathematical formulations reveal catastrophic points of failure when pushed beyond their operating boundaries. What happens if a designer attempts to lower latency by cutting window length in half? What happens if audio is sampled at studio rates ($48\text{ kHz}$)? What happens if dense matrix evaluation is chosen over sparse compression? In Section 1.5, we subject our theoretical models to four rigorous diagnostic stress-tests, examining where edge silicon breaks when pushed to its limits.

\end{document}
